% ============================================================================
% LANGENTWURF LE-023: Einkaufen im Internet
% Profil: S (Senioren 65+) | Tier: 2 (Standard)
% ============================================================================

\documentclass[11pt,a4paper]{article}
\usepackage[utf8]{inputenc}
\usepackage[T1]{fontenc}
\usepackage[ngerman]{babel}
\usepackage{geometry}
\usepackage{fancyhdr}
\usepackage{lastpage}
\usepackage{xcolor}
\usepackage{tcolorbox}
\usepackage{enumitem}
\usepackage{longtable}
\usepackage{hyperref}
\usepackage{amssymb}
\usepackage{eurosym}

\geometry{left=2.5cm, right=2.5cm, top=2.5cm, bottom=2.5cm}
\definecolor{fach}{HTML}{b35c00}
\definecolor{gray}{HTML}{666666}
\definecolor{successgreen}{HTML}{2a7a4b}
\definecolor{warnred}{HTML}{c62828}

\tcbuselibrary{skins,breakable}
\newtcolorbox{scorebox}{ colback=white, colframe=fach, fonttitle=\bfseries, title=Didaktik-Score, breakable }
\newtcolorbox{warnbox}[1][]{ colback=warnred!5, colframe=warnred, fonttitle=\bfseries, title=#1, breakable }
\newtcolorbox{tippbox}[1][]{ colback=successgreen!5, colframe=successgreen, fonttitle=\bfseries, title=#1, breakable }

\pagestyle{fancy}
\fancyhf{}
\fancyhead[L]{\small\textcolor{gray}{Digitale Grundlagen -- 023 Einkaufen im Internet}}
\fancyhead[R]{\small\textcolor{gray}{LE Tier 2 / Profil S}}
\fancyfoot[C]{\small\thepage{} / \pageref{LastPage}}
\renewcommand{\headrulewidth}{0.4pt}

\begin{document}

\begin{titlepage}
    \centering
    \vspace*{2cm}
    {\Large\textcolor{gray}{Digitale Grundlagen -- Senioren 65+}}\\[2cm]
    {\Huge\textcolor{fach}{\textbf{Langentwurf}}}\\[0.5cm]
    {\large Tier 2 | Profil S}\\[2cm]
    {\LARGE\textbf{023 -- Sicher einkaufen im Internet}}\\[0.5cm]
    {\large Modul F: Sicherheit \& Finanzen}\\[2cm]
    \begin{tabular}{rl}
        \textbf{Zeitrahmen:} & 60--75 Minuten \\
        \textbf{Vorkenntnisse:} & Online-Banking (022), Sicherheit (010) \\
    \end{tabular}
    \vfill
    {\normalsize Stand: \today}
\end{titlepage}

\tableofcontents
\newpage

\section{Bedingungsanalyse}

\subsection{Ausgangssituation}
\begin{itemize}
    \item Angst vor Betrug (``unseriöse Shops'')
    \item Unsicherheit bei Bezahlmethoden
    \item Angst, persönliche Daten preiszugeben
    \item Wunsch nach Bequemlichkeit (Lieferung nach Hause)
    \item Erfahrung mit Fake-Shops fehlt
\end{itemize}

\subsection{Typische Probleme}
\begin{itemize}
    \item ``Ist dieser Shop vertrauenswürdig?''
    \item ``Wie bezahle ich am sichersten?''
    \item ``Was, wenn die Ware nicht kommt?''
    \item ``Kann ich zurückgeben?''
    \item ``Warum brauchen die meine Adresse?''
\end{itemize}

\section{Sachanalyse}

\subsection{Kernkonzepte}
\begin{enumerate}
    \item \textbf{Seriöse Shops:} Impressum, https, bekannte Namen
    \item \textbf{Bezahlmethoden:} PayPal, Rechnung, Kreditkarte, Vorkasse
    \item \textbf{Widerrufsrecht:} 14 Tage Rückgabe ohne Grund
    \item \textbf{Fake-Shops:} Betrügerische Seiten, die Geld kassieren
    \item \textbf{Käuferschutz:} Schutz bei PayPal und Kreditkarte
\end{enumerate}

\subsection{Alltagsanalogien}
\begin{tabular}{|l|p{8cm}|}
\hline
\textbf{Begriff} & \textbf{Analogie} \\
\hline
Seriöser Shop & Bekanntes Geschäft in der Fußgängerzone \\
Fake-Shop & Kofferraum-Verkäufer auf dem Parkplatz \\
Warenkorb & Einkaufswagen im Supermarkt \\
Checkout & Kasse im Geschäft \\
Widerrufsrecht & Umtausch im Laden (aber besser: 14 Tage!) \\
\hline
\end{tabular}

\section{Didaktische Analyse}

\subsection{Stellung in der Reihe}
Ergänzt Online-Banking (022). Setzt Sicherheitsgrundlagen voraus.

\subsection{Didaktische Reduktion}
\textbf{Ausgeklammert:} Internationale Bestellungen, Zoll, Marketplace-Verkäufer, Abo-Fallen im Detail.

\textbf{Fokus:} Seriöse Shops erkennen, sichere Bezahlmethoden wählen, Rechte kennen.

\section{Lernziele}

\begin{tabular}{|c|p{7cm}|c|c|}
\hline
\textbf{Nr.} & \textbf{Die Teilnehmer können...} & \textbf{Stufe} & \textbf{Niveau} \\
\hline
FZ1 & 3 Merkmale seriöser Shops nennen & Wissen & Basis \\
FZ2 & Fake-Shops erkennen (Warnsignale) & Analysieren & Basis \\
FZ3 & sichere Bezahlmethoden benennen & Wissen & Basis \\
FZ4 & ihr Widerrufsrecht erklären & Verstehen & Basis \\
FZ5 & einen Bestellvorgang nachvollziehen & Verstehen & Vertiefung \\
\hline
\end{tabular}

\section{Methodische Analyse}

\begin{tippbox}[3 Merkmale seriöser Shops]
\begin{enumerate}
    \item \textbf{https://} und Schloss-Symbol in der Adresszeile
    \item \textbf{Impressum} mit vollständiger Adresse und Telefonnummer
    \item \textbf{Bekannte Namen:} Amazon, Otto, MediaMarkt, etc.
\end{enumerate}
\end{tippbox}

\begin{warnbox}[Fake-Shop Warnsignale]
\begin{itemize}
    \item Preise \textbf{viel zu günstig} (zu schön, um wahr zu sein)
    \item \textbf{Nur Vorkasse} als Zahlungsmethode
    \item \textbf{Kein Impressum} oder ausländische Adresse
    \item \textbf{Schlechtes Deutsch} (Übersetzungsfehler)
    \item \textbf{Keine Bewertungen} oder nur 5-Sterne-Bewertungen
\end{itemize}
\end{warnbox}

\subsection{Bezahlmethoden im Vergleich}
\begin{tabular}{|l|c|c|p{4cm}|}
\hline
\textbf{Methode} & \textbf{Sicherheit} & \textbf{Empfehlung} & \textbf{Hinweis} \\
\hline
Rechnung & Sehr hoch & \textcolor{successgreen}{\textbf{Ja}} & Erst zahlen nach Erhalt \\
PayPal & Hoch & \textcolor{successgreen}{\textbf{Ja}} & Käuferschutz inklusive \\
Kreditkarte & Hoch & Ja & Rückbuchung möglich \\
Lastschrift & Mittel & Bedingt & Rückbuchung 8 Wochen \\
Vorkasse & Niedrig & \textcolor{warnred}{\textbf{Nein}} & Geld weg bei Betrug! \\
\hline
\end{tabular}

\section{Verlaufsplanung}

\begin{longtable}{|p{1cm}|p{2cm}|p{5cm}|p{3cm}|p{2cm}|}
\hline
\textbf{Zeit} & \textbf{Phase} & \textbf{Inhalt} & \textbf{TN-Aktivität} & \textbf{Material} \\
\hline
0--10 & Einstieg & ``Haben Sie schon mal online bestellt? Was war gut/schlecht?'' & Berichten & -- \\
\hline
10--25 & Seriöse Shops & 3 Merkmale vorstellen, Beispiele zeigen (Amazon, Otto) & Mitschreiben & S.1 \\
\hline
25--40 & Fake-Shops & Warnsignale zeigen, echte Beispiele (Screenshots) & Analysieren & S.1 \\
\hline
40--55 & Bezahlen & Methoden vergleichen, PayPal/Rechnung empfehlen & Nachvollziehen & S.2 \\
\hline
55--70 & Rechte & Widerrufsrecht, Rückgabe, Verbraucherzentrale & Notieren & S.2-3 \\
\hline
\end{longtable}

\section{Lösungen}

\textbf{3 Merkmale seriöser Shops:}
\begin{enumerate}
    \item https:// und Schloss-Symbol
    \item Impressum mit vollständiger Adresse
    \item Bekannte, etablierte Shops
\end{enumerate}

\textbf{Widerrufsrecht:}
\begin{itemize}
    \item 14 Tage Zeit zum Widerruf (ohne Begründung!)
    \item Ab Erhalt der Ware
    \item Rücksendung auf eigene Kosten (meist)
    \item Geld zurück innerhalb 14 Tagen nach Widerruf
\end{itemize}

\textbf{Bei Problemen:}
\begin{itemize}
    \item Zuerst: Shop kontaktieren (schriftlich!)
    \item Bei PayPal: Käuferschutz-Fall öffnen
    \item Verbraucherzentrale: Beratung und Hilfe
    \item Polizei: Bei Betrug Anzeige erstatten
\end{itemize}

\textbf{Empfohlene Shops für Anfänger:}
\begin{itemize}
    \item \textbf{Amazon:} Größte Auswahl, einfache Rückgabe
    \item \textbf{Otto:} Deutscher Versandhändler, Kauf auf Rechnung
    \item \textbf{MediaMarkt/Saturn:} Elektronik
    \item \textbf{dm/Rossmann:} Drogerie
    \item \textbf{Shops lokaler Händler:} Unterstützt die Region
\end{itemize}

\section{Didaktik-Score}

\begin{scorebox}
\begin{tabular}{|c|l|c|c|p{3.5cm}|}
\hline
\# & Dimension & Gew. & Score & Notiz \\
\hline
1 & Differenzierung & 15\% & 8/10 & Basis/Vertiefung ok \\
2 & Taxonomie & 10\% & 9/10 & Wissen + Analysieren \\
3 & Alltagsrelevanz & 10\% & 10/10 & Sehr hohe Relevanz \\
4 & Aktivierung & 10\% & 8/10 & Analyse von Beispielen \\
5 & Scaffolding & 10\% & 9/10 & Klare Checklisten \\
6 & Feedback & 10\% & 8/10 & Ermutigung wichtig \\
7 & Sprachliche Klarheit & 10\% & 9/10 & Einfach erklärt \\
8 & Motivation & 10\% & 10/10 & Praktischer Nutzen hoch \\
9 & Barrierefreiheit & 10\% & 9/10 & Große Schrift \\
10 & Praktikabilität & 5\% & 8/10 & 70 Min gut \\
\hline
\multicolumn{3}{|r|}{\textbf{GESAMT:}} & \textbf{9.05/10} & \textcolor{successgreen}{\textbf{FREIGABE}} \\
\hline
\end{tabular}
\end{scorebox}

\end{document}
